\subsection{Алгоритм Дейкстры}

\subsection*{Генерация графа}

Функция \texttt{generateGraph} генерирует случайный неориентированный граф с \texttt{n} вершинами и до \texttt{m} ребер. Вес ребер присваиваются случайным образом в диапазоне от \texttt{q} до \texttt{r}.

Максимальное количество ребер ограничено формулой \( \frac{n(n-1)}{2} \), которая представляет количество ребер в полном графе. Двумерный вектор используется для отслеживания существующих ребер, чтобы избежать дублирования. Ребра добавляются в список смежности до достижения желаемого количества ребер.

\begin{lstlisting}
    while edge_count < m do
    begin
      i := Random(n) + 1;
      j := Random(n) + 1;

      if (i <> j) and (not was[i][j]) then
      begin
        weight := Random(r - q + 1) + q;
        SetLength(adj_list[i], Length(adj_list[i]) + 1);
        adj_list[i][High(adj_list[i])].Node := j;
        adj_list[i][High(adj_list[i])].Weight := weight;

        SetLength(adj_list[j], Length(adj_list[j]) + 1);
        adj_list[j][High(adj_list[j])].Node := i;
        adj_list[j][High(adj_list[j])].Weight := weight;

        was[i][j] := True;
        was[j][i] := True;
        Inc(edge_count);
      end;
    end;
\end{lstlisting}

\textbf{Анализ сложности:}
\begin{itemize}
    \item Генерация случайных индексов \(i\) и \(j\) для вершины выполняется за \(O(1)\).
    \item Проверка уникальности ребра через массив \texttt{was} занимает \(O(1)\).
    \item Обновление списка смежности требует \(O(1)\) для каждого добавления ребра.
\end{itemize}

Процесс продолжается, пока не будет создано \(m\) уникальных рёбер. Для полного графа (\(m = \frac{n(n-1)}{2}\)) сложность составляет \(O(m)\). С учётом возможных повторов, средняя сложность генерации графа остаётся \(O(m)\).

\subsection*{Алгоритм Дейкстры}

Функция реализует алгоритм Дейкстры для нахождения кратчайшего пути от стартовой вершины ко всем другим вершинам графа.

Алгоритм инициализирует расстояния до бесконечности (\texttt{INT\_MAX}) и устанавливает расстояние до стартовой вершины равным нулю. Приоритетная очередь, представляемая (\texttt{DHeap}), используется для эффективного извлечения вершины с минимальным расстоянием. 

Для каждой соседней вершины проверяется, можно ли улучшить текущее известное расстояние до неё. Если новое расстояние, вычисленное как сумма расстояния до текущей вершины и веса ребра, ведущего к соседней вершине, меньше, чем уже известное расстояние, то оно обновляется. Также в таком случае в кучу добавляется информация о новой вершине с обновленным расстоянием.

Процесс продолжается, пока не будут обработаны все доступные вершины. В итоге алгоритм возвращает массив расстояний, в котором содержатся кратчайшие расстояния от стартовой вершины до всех остальных вершин графа.

\begin{lstlisting}
    while not dHeap.IsEmpty() do
    begin
        minPair := dHeap.ExtractMin();
        i := minPair.first;
        currentDist := minPair.second;

        if processed[i] then 
            Continue; 

        processed[i] := True;

        for k := 0 to High(graph[i]) do
        begin
            j := graph[i][k].first;
            newDist := graph[i][k].second + currentDist;

            if newDist < dist[j] then
            begin
                dist[j] := newDist;
                up[j] := i;
                dHeap.Insert(j, dist[j]);
            end;
        end;
    end;
\end{lstlisting}

\textbf{Анализ сложности:}
\begin{itemize}
    \item Инициализация массива расстояний: \(O(n)\).
    \item Извлечение вершины из очереди: \(O(\log n)\) на операцию.
    \item Обработка всех соседей вершины: \(O(\text{степень вершины})\).
    \item Обновление расстояний для всех \(m\) рёбер требует \(O(m \log n)\).
\end{itemize}

Итоговая сложность алгоритма Дейкстры:
\[
O(n \log n + m \log n) = O((n + m) \log n),
\]
где \(n\) — количество вершин, а \(m\) — количество рёбер.
